% Input representation: identical for every task, so each model paper inputs
% this file under its own heading and then adds what is specific to its task.
% The prose follows docs/writing.md.
All tasks in this project share one input representation, so two tasks differ
in their target and in the information a model is given, and not in how a
position is encoded.

\paragraph{Square tokens.}
A position is encoded as 64 tokens, one per square. Token 0 denotes an empty
square, tokens 1 to 6 the pawn, knight, bishop, rook, queen and king of the
side to move, which we call the \emph{mover}, and tokens 7 to 12 the
opponent's pieces. When Black is to move, the board is flipped vertically and
the colours are swapped, so that the mover always plays up the board. A
position with White to move and the corresponding colour-reversed position
with Black to move therefore have the same encoding, and a model does not
have to learn the two cases separately.

\paragraph{Auxiliary input planes.}
Information about individual squares is supplied as binary planes over the 64
squares, which a network receives together with the tokens. This task uses
\V{task.marks.count} planes (Table~\ref{tab:marks}). The first marks the en
passant target square, and the others mark the origin and destination squares
of moves. A plane identifies the squares of a move and carries no other
information about it, such as whether the move is a capture or a check.

\begin{table}[htbp]\centering\small
\block{marks}
\caption{Auxiliary input planes. Squares are given from the mover's point of
view, after the board has been flipped.}
\label{tab:marks}
\end{table}

\ifsequence
\paragraph{Sequences of positions.}
An example of this task is not a single position. It is the sequence of one
player's moves in a game, and each element of the sequence is a position
encoded as above: the board before the move, with the opponent's preceding
move and the move played marked in the auxiliary planes. Scalar inputs
therefore come at two levels, one for each move and one for the game.

\paragraph{Per-move features.}
Each position of the sequence has \V{task.steps.count} scalar features of its
own (Table~\ref{tab:steps}): the four castling rights and the task's
information about that move. A network receives them together with the
position they belong to.

\begin{table}[htbp]\centering\small
\block{steps}
\caption{Per-move features: the scalar inputs attached to each position of a
sequence.}
\label{tab:steps}
\end{table}

\paragraph{Scalar context features.}
Information about the game as a whole is given as \V{task.context.count}
scalar context features (Table~\ref{tab:reads}). Every model receives them,
including the gradient-boosted trees, which never see a board. For a network,
both kinds of scalar feature are standardised with the mean and standard
deviation of the training set.

\begin{table}[htbp]\centering\small
\block{reads}
\caption{Scalar context features: the game-level inputs given to every
model.}
\label{tab:reads}
\end{table}

\paragraph{Hand-crafted features.}
Gradient-boosted trees cannot take a sequence of boards as input, so we also
compute \V{task.hand.count} aggregate features over the player's moves
(Table~\ref{tab:hand} in the appendix): how much win probability the moves
lost on average, at worst and how often beyond fixed thresholds, how the clock
was used, and how many moves were available. These features have two roles.
They are the inputs of the GBDT baseline, and so define the level that a
network reading every position has to exceed. They can also be given to a
network as additional context features, which tests whether the network
already extracts the same information from the sequence.

\else
\paragraph{Scalar context features.}
Information that does not belong to a square is given as scalar context
features: the four castling rights and the scalar inputs specific to the
task. This task uses \V{task.context.count} (Table~\ref{tab:reads}). Every
model receives them, including the gradient-boosted trees, which never see
the board. For a network they are standardised with the mean and standard
deviation of the training set.

\begin{table}[htbp]\centering\small
\block{reads}
\caption{Scalar context features: the inputs given to every model in addition
to the board.}
\label{tab:reads}
\end{table}

\paragraph{Hand-crafted features.}
Gradient-boosted trees cannot take 64 tokens as input, so we also compute
\V{task.hand.count} hand-crafted features of the position
(Table~\ref{tab:hand} in the appendix): material, mobility, available checks
and captures, hanging pieces, passed pawns and attacks near each king, and,
where the task gives a line of moves, features of that line. These features
have two roles. They are the inputs of the GBDT baseline, and so define the
level that a network reading the board has to exceed. They can also be given
to a network as additional context features, which tests whether the network
already extracts the same information from the board.

\fi

\iflook
\begin{figure}[htbp]\centering
\block{encoding}
\ifsequence
\caption{One position of one example, shown twice: the position before the
player's most costly move of the game. Left: as played, with arrows for the
opponent's preceding move (grey) and the player's move (blue). Right: as
encoded, oriented so that the mover plays up the board, with the squares set
in an auxiliary plane shaded grey for the opponent's move and blue for the
player's move. Every position of the sequence is encoded in this way.}
\else
\caption{One example from the dataset, shown twice. Left: as played, with
arrows for the opponent's preceding move (grey) and the move the example
concerns (blue). Right: as encoded, oriented so that the mover plays up the
board, with the squares set in an auxiliary plane shaded grey for the
opponent's move and, where the task marks a line of moves, blue for the
mover's moves and orange for the replies.}
\fi
\label{fig:encoding}
\end{figure}
\fi
